\chapter{From Signal to Vector}
\label{ch:signal-to-vector}

Raw observations arrive as pixels, tokens, or sensor readings. The first modelling choice is the space in which those observations become comparable.

A vector $x \in \R^d$ is not merely a container. Its dimension bounds capacity; an inner product defines similarity; a basis decides which variations count as independent. The geometry of that space is already a theory of what can be represented.

\begin{center}
\begin{tikzpicture}[scale=1.1]
  \draw[->] (0,0) -- (3.2,0) node[right] {$x$};
  \draw[->] (0,0) -- (0,2.6) node[above] {$y$};
  \draw[thick,accent] (0,0) -- (2.4,1.5) node[above right] {$u$};
  \draw[thick] (0,0) -- (2.1,0) node[below] {$v$};
  \draw[dotted] (2.4,1.5) -- (2.1,0);
\end{tikzpicture}

\end{center}

\begin{remark}
Later chapters return to the same geometry when attention, retrieval, and contrastive learning reuse inner products as scores.
\end{remark}
